% PKG-5 "Explorable textbook: one scene, three views" -- one-page print specimen (design-packages-B, 2026-10-09)
% Fonts: Libertinus Serif / Math / Sans / Mono (OFL). Colours: ../tokens/pkg5-tokens.sty (generated).
% The figure is \input from scene/scene_2pz.tikz, emitted by scene/scene_2pz.py from the scene graph.
\documentclass[fontsize=11pt,parskip=false]{scrartcl}
\usepackage[a4paper,inner=24mm,textwidth=112mm,marginparsep=6mm,marginparwidth=44mm,
  top=22mm,bottom=26mm]{geometry}
\usepackage{fontspec}
\usepackage{amsmath}
\usepackage[math-style=ISO,bold-style=ISO]{unicode-math}
\setmainfont{LibertinusSerif}[Extension=.otf,UprightFont=*-Regular,ItalicFont=*-Italic,
  BoldFont=*-Semibold,BoldItalicFont=*-SemiboldItalic]
\setsansfont{LibertinusSans}[Extension=.otf,UprightFont=*-Regular,ItalicFont=*-Italic,BoldFont=*-Bold]
\setmonofont{LibertinusMono-Regular.otf}
\setmathfont{LibertinusMath-Regular.otf}
\usepackage[ngerman]{babel}
\usepackage[autostyle]{csquotes}
\usepackage[protrusion=true,expansion=false]{microtype}
\usepackage[locale=DE]{siunitx}
\usepackage{xcolor}
% pkg5-tokens.sty -- GENERATED by build_tokens.py from pkg5.tokens.json. Do not edit.
\definecolor{tpc-paper}{HTML}{FFFFFF}
\definecolor{tpcvid-paper}{HTML}{FAF9F6}
\definecolor{tpc-plate}{HTML}{FFFFFF}
\definecolor{tpcvid-plate}{HTML}{FAF9F6}
\definecolor{tpc-ink}{HTML}{161616}
\definecolor{tpcvid-ink}{HTML}{161616}
\definecolor{tpc-ink-muted}{HTML}{5F5D5A}
\definecolor{tpcvid-ink-muted}{HTML}{5F5D5A}
\definecolor{tpc-rule}{HTML}{C5C4C1}
\definecolor{tpcvid-rule}{HTML}{C5C4C1}
\definecolor{tpc-accent}{HTML}{161616}
\definecolor{tpcvid-accent}{HTML}{161616}
\definecolor{tpc-link}{HTML}{302E28}
\definecolor{tpcvid-link}{HTML}{302E28}
\definecolor{tpc-state}{HTML}{014069}
\definecolor{tpcvid-state}{HTML}{014069}
\definecolor{tpc-effect}{HTML}{017A58}
\definecolor{tpcvid-effect}{HTML}{017A58}
\definecolor{tpc-dynamics}{HTML}{D6781C}
\definecolor{tpcvid-dynamics}{HTML}{D6781C}
\definecolor{tpc-environment}{HTML}{CCCAC6}
\definecolor{tpcvid-environment}{HTML}{CCCAC6}

\usepackage{tikz}
\usetikzlibrary{arrows.meta,calc}
\usepackage[most]{tcolorbox}
\color{tpc-ink}
\linespread{1.06}  % 11 on about 14.5
\newcommand{\dd}{\mathrm{d}}\newcommand{\eu}{\mathrm{e}}\newcommand{\iu}{\mathrm{i}}
\newcommand{\tr}{\operatorname{tr}}
\newcommand{\side}[1]{\marginpar{\raggedright\fontsize{8.5}{11}\selectfont\color{tpc-ink-muted}#1}}
\newcommand{\live}[1]{\marginpar{\raggedright\sffamily\fontsize{8}{10.5}\selectfont
  \tikz[baseline=-0.6ex]\fill[tpc-ink](0,-0.9mm)--(1.6mm,0)--(0,0.9mm)--cycle;\ \textbf{Live}\enspace#1}}
% theorem: hairlines above and below, small-caps label, no fill (colour is reserved for physics)
\newtcolorbox{postulat}[1]{enhanced,breakable,colback=tpc-paper,colframe=tpc-paper,boxrule=0pt,arc=0pt,
  left=0pt,right=0pt,top=1.5mm,bottom=1.5mm,
  borderline north={0.4pt}{0pt}{tpc-ink},borderline south={0.4pt}{0pt}{tpc-ink},
  before upper={\noindent{\scshape #1.}\enskip}}
% Vertiefung (depth) block, after layered-docs/REPORT.md: 9.5 pt, 0.4 pt outer hairline, run-in small-caps label, tag in the margin
\newtcolorbox{zoomin}[2]{enhanced,breakable,colback=tpc-paper,colframe=tpc-paper,boxrule=0pt,arc=0pt,
  left=0pt,right=0pt,top=0.5mm,bottom=0.5mm,fontupper=\fontsize{9.5}{12.5}\selectfont,
  overlay={\draw[tpc-ink,line width=0.4pt]([xshift=3mm]frame.north east)--([xshift=3mm]frame.south east);
    \node[anchor=north west,inner sep=0pt] at ([xshift=6mm]frame.north east)
      {\tikz[baseline]\node[draw=tpc-ink,circle,line width=0.4pt,inner sep=0.8pt,font=\sffamily\scriptsize\bfseries]{V};%
       \ \sffamily\fontsize{8}{10}\selectfont\color{tpc-ink-muted}#2};}}
\pagestyle{empty}
\begin{document}
\noindent{\sffamily\fontsize{8.5}{10}\selectfont\color{tpc-ink-muted}16\enspace Das Wasserstoffatom}\\[1.5mm]
{\fontsize{20}{24}\selectfont Wo ist das Elektron?}
\vspace{5mm}

\noindent Der Zustand \(2\mathrm{p}_z\) des Wasserstoffatoms ist die Eigenfunktion zu \(n=2\), \(\ell=1\), \(m=0\).
In atomaren Einheiten (\(a_0=1\)) lautet sie
\begin{equation}
  \psi_{210}(r,\vartheta,\varphi) = \frac{1}{4\sqrt{2\pi}}\; r\,\eu^{-r/2}\cos\vartheta ,
  \qquad
  \lvert\psi_{210}\rvert^2 = \frac{r^2\,\eu^{-r}\cos^2\vartheta}{32\pi}.
\end{equation}
Die Dichte hängt nicht von \(\varphi\) ab, verschwindet in der Ebene \(z=0\) und ist auf der
\(z\)-Achse bei \(r=2\) maximal.%
\side{\(\vartheta\) ist der Winkel zur \(z\)-Achse, nicht zur Ebene. Die Knotenebene ist eine
Eigenschaft des \emph{Zustands}, darum trägt sie dessen Farbe.}

\begin{postulat}{Postulat P4 (Bornsche Regel), Ortsform}
Ist das System im Zustand \(\rho=\lvert\psi\rangle\langle\psi\rvert\) und spricht der Detektor auf das Gebiet
\(V\subset\mathbb{R}^3\) an, so ist \(E_V=\chi_V(\hat{\mathbf{x}})\) der Effekt und
\begin{equation}
  p(V) = \tr(\rho\,E_V) = \int_V \lvert\psi(\mathbf{x})\rvert^2\,\dd^3x .
\end{equation}
\end{postulat}

\begin{center}
% GENERATED by scene_2pz.py from the scene graph; do not edit.
\begin{tikzpicture}[font=\sffamily\footnotesize]
\fill[tpc-state,opacity=0.18] (4.87mm,-5.60mm) -- (5.89mm,-7.29mm) -- (6.21mm,-8.27mm) -- (6.40mm,-9.15mm) -- (6.50mm,-9.97mm) -- (6.53mm,-10.76mm) -- (6.52mm,-11.52mm) -- (6.47mm,-12.26mm) -- (6.38mm,-12.99mm) -- (6.27mm,-13.69mm) -- (6.12mm,-14.38mm) -- (5.95mm,-15.05mm) -- (5.76mm,-15.70mm) -- (5.54mm,-16.33mm) -- (5.31mm,-16.94mm) -- (5.06mm,-17.53mm) -- (4.79mm,-18.10mm) -- (4.51mm,-18.64mm) -- (4.21mm,-19.17mm) -- (3.90mm,-19.67mm) -- (3.58mm,-20.14mm) -- (3.25mm,-20.60mm) -- (2.91mm,-21.02mm) -- (2.57mm,-21.42mm) -- (2.21mm,-21.80mm) -- (1.85mm,-22.14mm) -- (1.49mm,-22.46mm) -- (1.12mm,-22.75mm) -- (0.75mm,-23.00mm) -- (0.37mm,-23.23mm) -- (0.00mm,-23.43mm) -- (-0.37mm,-23.60mm) -- (-0.75mm,-23.73mm) -- (-1.12mm,-23.84mm) -- (-1.49mm,-23.91mm) -- (-1.85mm,-23.95mm) -- (-2.21mm,-23.96mm) -- (-2.57mm,-23.93mm) -- (-2.91mm,-23.87mm) -- (-3.25mm,-23.77mm) -- (-3.58mm,-23.64mm) -- (-3.90mm,-23.48mm) -- (-4.21mm,-23.28mm) -- (-4.51mm,-23.04mm) -- (-4.79mm,-22.77mm) -- (-5.06mm,-22.47mm) -- (-5.31mm,-22.12mm) -- (-5.54mm,-21.74mm) -- (-5.76mm,-21.32mm) -- (-5.95mm,-20.86mm) -- (-6.12mm,-20.36mm) -- (-6.27mm,-19.81mm) -- (-6.38mm,-19.22mm) -- (-6.47mm,-18.59mm) -- (-6.52mm,-17.89mm) -- (-6.53mm,-17.14mm) -- (-6.50mm,-16.32mm) -- (-6.40mm,-15.40mm) -- (-6.21mm,-14.34mm) -- (-5.89mm,-13.05mm) -- (-4.87mm,-10.35mm) -- (-4.87mm,-10.35mm) -- (-3.75mm,-8.30mm) -- (-3.30mm,-7.62mm) -- (-2.98mm,-7.16mm) -- (-2.71mm,-6.81mm) -- (-2.49mm,-6.53mm) -- (-2.29mm,-6.29mm) -- (-2.12mm,-6.09mm) -- (-1.96mm,-5.91mm) -- (-1.82mm,-5.75mm) -- (-1.69mm,-5.62mm) -- (-1.57mm,-5.49mm) -- (-1.45mm,-5.38mm) -- (-1.34mm,-5.27mm) -- (-1.24mm,-5.18mm) -- (-1.15mm,-5.09mm) -- (-1.05mm,-5.01mm) -- (-0.96mm,-4.94mm) -- (-0.88mm,-4.87mm) -- (-0.80mm,-4.80mm) -- (-0.72mm,-4.74mm) -- (-0.64mm,-4.68mm) -- (-0.56mm,-4.63mm) -- (-0.49mm,-4.58mm) -- (-0.42mm,-4.53mm) -- (-0.35mm,-4.48mm) -- (-0.28mm,-4.44mm) -- (-0.21mm,-4.40mm) -- (-0.14mm,-4.36mm) -- (-0.07mm,-4.33mm) -- (0.00mm,-4.29mm) -- (0.07mm,-4.26mm) -- (0.14mm,-4.23mm) -- (0.21mm,-4.20mm) -- (0.28mm,-4.17mm) -- (0.35mm,-4.15mm) -- (0.42mm,-4.12mm) -- (0.49mm,-4.10mm) -- (0.56mm,-4.08mm) -- (0.64mm,-4.06mm) -- (0.72mm,-4.04mm) -- (0.80mm,-4.02mm) -- (0.88mm,-4.01mm) -- (0.96mm,-3.99mm) -- (1.05mm,-3.98mm) -- (1.15mm,-3.97mm) -- (1.24mm,-3.97mm) -- (1.34mm,-3.96mm) -- (1.45mm,-3.96mm) -- (1.57mm,-3.96mm) -- (1.69mm,-3.97mm) -- (1.82mm,-3.98mm) -- (1.96mm,-3.99mm) -- (2.12mm,-4.02mm) -- (2.29mm,-4.05mm) -- (2.49mm,-4.10mm) -- (2.71mm,-4.16mm) -- (2.98mm,-4.25mm) -- (3.30mm,-4.39mm) -- (3.75mm,-4.64mm) -- (4.87mm,-5.60mm) -- cycle;
\draw[tpc-state,line width=1.0pt] (4.87mm,-5.60mm) -- (5.89mm,-7.29mm) -- (6.21mm,-8.27mm) -- (6.40mm,-9.15mm) -- (6.50mm,-9.97mm) -- (6.53mm,-10.76mm) -- (6.52mm,-11.52mm) -- (6.47mm,-12.26mm) -- (6.38mm,-12.99mm) -- (6.27mm,-13.69mm) -- (6.12mm,-14.38mm) -- (5.95mm,-15.05mm) -- (5.76mm,-15.70mm) -- (5.54mm,-16.33mm) -- (5.31mm,-16.94mm) -- (5.06mm,-17.53mm) -- (4.79mm,-18.10mm) -- (4.51mm,-18.64mm) -- (4.21mm,-19.17mm) -- (3.90mm,-19.67mm) -- (3.58mm,-20.14mm) -- (3.25mm,-20.60mm) -- (2.91mm,-21.02mm) -- (2.57mm,-21.42mm) -- (2.21mm,-21.80mm) -- (1.85mm,-22.14mm) -- (1.49mm,-22.46mm) -- (1.12mm,-22.75mm) -- (0.75mm,-23.00mm) -- (0.37mm,-23.23mm) -- (0.00mm,-23.43mm) -- (-0.37mm,-23.60mm) -- (-0.75mm,-23.73mm) -- (-1.12mm,-23.84mm) -- (-1.49mm,-23.91mm) -- (-1.85mm,-23.95mm) -- (-2.21mm,-23.96mm) -- (-2.57mm,-23.93mm) -- (-2.91mm,-23.87mm) -- (-3.25mm,-23.77mm) -- (-3.58mm,-23.64mm) -- (-3.90mm,-23.48mm) -- (-4.21mm,-23.28mm) -- (-4.51mm,-23.04mm) -- (-4.79mm,-22.77mm) -- (-5.06mm,-22.47mm) -- (-5.31mm,-22.12mm) -- (-5.54mm,-21.74mm) -- (-5.76mm,-21.32mm) -- (-5.95mm,-20.86mm) -- (-6.12mm,-20.36mm) -- (-6.27mm,-19.81mm) -- (-6.38mm,-19.22mm) -- (-6.47mm,-18.59mm) -- (-6.52mm,-17.89mm) -- (-6.53mm,-17.14mm) -- (-6.50mm,-16.32mm) -- (-6.40mm,-15.40mm) -- (-6.21mm,-14.34mm) -- (-5.89mm,-13.05mm) -- (-4.87mm,-10.35mm) -- (-4.87mm,-10.35mm) -- (-3.75mm,-8.30mm) -- (-3.30mm,-7.62mm) -- (-2.98mm,-7.16mm) -- (-2.71mm,-6.81mm) -- (-2.49mm,-6.53mm) -- (-2.29mm,-6.29mm) -- (-2.12mm,-6.09mm) -- (-1.96mm,-5.91mm) -- (-1.82mm,-5.75mm) -- (-1.69mm,-5.62mm) -- (-1.57mm,-5.49mm) -- (-1.45mm,-5.38mm) -- (-1.34mm,-5.27mm) -- (-1.24mm,-5.18mm) -- (-1.15mm,-5.09mm) -- (-1.05mm,-5.01mm) -- (-0.96mm,-4.94mm) -- (-0.88mm,-4.87mm) -- (-0.80mm,-4.80mm) -- (-0.72mm,-4.74mm) -- (-0.64mm,-4.68mm) -- (-0.56mm,-4.63mm) -- (-0.49mm,-4.58mm) -- (-0.42mm,-4.53mm) -- (-0.35mm,-4.48mm) -- (-0.28mm,-4.44mm) -- (-0.21mm,-4.40mm) -- (-0.14mm,-4.36mm) -- (-0.07mm,-4.33mm) -- (0.00mm,-4.29mm) -- (0.07mm,-4.26mm) -- (0.14mm,-4.23mm) -- (0.21mm,-4.20mm) -- (0.28mm,-4.17mm) -- (0.35mm,-4.15mm) -- (0.42mm,-4.12mm) -- (0.49mm,-4.10mm) -- (0.56mm,-4.08mm) -- (0.64mm,-4.06mm) -- (0.72mm,-4.04mm) -- (0.80mm,-4.02mm) -- (0.88mm,-4.01mm) -- (0.96mm,-3.99mm) -- (1.05mm,-3.98mm) -- (1.15mm,-3.97mm) -- (1.24mm,-3.97mm) -- (1.34mm,-3.96mm) -- (1.45mm,-3.96mm) -- (1.57mm,-3.96mm) -- (1.69mm,-3.97mm) -- (1.82mm,-3.98mm) -- (1.96mm,-3.99mm) -- (2.12mm,-4.02mm) -- (2.29mm,-4.05mm) -- (2.49mm,-4.10mm) -- (2.71mm,-4.16mm) -- (2.98mm,-4.25mm) -- (3.30mm,-4.39mm) -- (3.75mm,-4.64mm) -- (4.87mm,-5.60mm) -- cycle;
\fill[tpc-state,opacity=0.06] (-5.30mm,10.29mm) -- (-30.08mm,-1.81mm) -- (5.30mm,-10.29mm) -- (30.08mm,1.81mm) -- cycle;
\draw[tpc-state,line width=0.5pt,dash pattern=on 2pt off 1.5pt] (-5.30mm,10.29mm) -- (-30.08mm,-1.81mm) -- (5.30mm,-10.29mm) -- (30.08mm,1.81mm) -- cycle;
\draw[tpc-ink,line width=0.4pt] (0.00mm,-29.32mm) -- (0.00mm,29.32mm);
\fill[tpc-state,opacity=0.18] (4.87mm,10.35mm) -- (5.89mm,13.05mm) -- (6.21mm,14.34mm) -- (6.40mm,15.40mm) -- (6.50mm,16.32mm) -- (6.53mm,17.14mm) -- (6.52mm,17.89mm) -- (6.47mm,18.59mm) -- (6.38mm,19.22mm) -- (6.27mm,19.81mm) -- (6.12mm,20.36mm) -- (5.95mm,20.86mm) -- (5.76mm,21.32mm) -- (5.54mm,21.74mm) -- (5.31mm,22.12mm) -- (5.06mm,22.47mm) -- (4.79mm,22.77mm) -- (4.51mm,23.04mm) -- (4.21mm,23.28mm) -- (3.90mm,23.48mm) -- (3.58mm,23.64mm) -- (3.25mm,23.77mm) -- (2.91mm,23.87mm) -- (2.57mm,23.93mm) -- (2.21mm,23.96mm) -- (1.85mm,23.95mm) -- (1.49mm,23.91mm) -- (1.12mm,23.84mm) -- (0.75mm,23.73mm) -- (0.37mm,23.60mm) -- (0.00mm,23.43mm) -- (-0.37mm,23.23mm) -- (-0.75mm,23.00mm) -- (-1.12mm,22.75mm) -- (-1.49mm,22.46mm) -- (-1.85mm,22.14mm) -- (-2.21mm,21.80mm) -- (-2.57mm,21.42mm) -- (-2.91mm,21.02mm) -- (-3.25mm,20.60mm) -- (-3.58mm,20.14mm) -- (-3.90mm,19.67mm) -- (-4.21mm,19.17mm) -- (-4.51mm,18.64mm) -- (-4.79mm,18.10mm) -- (-5.06mm,17.53mm) -- (-5.31mm,16.94mm) -- (-5.54mm,16.33mm) -- (-5.76mm,15.70mm) -- (-5.95mm,15.05mm) -- (-6.12mm,14.38mm) -- (-6.27mm,13.69mm) -- (-6.38mm,12.99mm) -- (-6.47mm,12.26mm) -- (-6.52mm,11.52mm) -- (-6.53mm,10.76mm) -- (-6.50mm,9.97mm) -- (-6.40mm,9.15mm) -- (-6.21mm,8.27mm) -- (-5.89mm,7.29mm) -- (-4.87mm,5.60mm) -- (-4.87mm,5.60mm) -- (-3.75mm,4.64mm) -- (-3.30mm,4.39mm) -- (-2.98mm,4.25mm) -- (-2.71mm,4.16mm) -- (-2.49mm,4.10mm) -- (-2.29mm,4.05mm) -- (-2.12mm,4.02mm) -- (-1.96mm,3.99mm) -- (-1.82mm,3.98mm) -- (-1.69mm,3.97mm) -- (-1.57mm,3.96mm) -- (-1.45mm,3.96mm) -- (-1.34mm,3.96mm) -- (-1.24mm,3.97mm) -- (-1.15mm,3.97mm) -- (-1.05mm,3.98mm) -- (-0.96mm,3.99mm) -- (-0.88mm,4.01mm) -- (-0.80mm,4.02mm) -- (-0.72mm,4.04mm) -- (-0.64mm,4.06mm) -- (-0.56mm,4.08mm) -- (-0.49mm,4.10mm) -- (-0.42mm,4.12mm) -- (-0.35mm,4.15mm) -- (-0.28mm,4.17mm) -- (-0.21mm,4.20mm) -- (-0.14mm,4.23mm) -- (-0.07mm,4.26mm) -- (0.00mm,4.29mm) -- (0.07mm,4.33mm) -- (0.14mm,4.36mm) -- (0.21mm,4.40mm) -- (0.28mm,4.44mm) -- (0.35mm,4.48mm) -- (0.42mm,4.53mm) -- (0.49mm,4.58mm) -- (0.56mm,4.63mm) -- (0.64mm,4.68mm) -- (0.72mm,4.74mm) -- (0.80mm,4.80mm) -- (0.88mm,4.87mm) -- (0.96mm,4.94mm) -- (1.05mm,5.01mm) -- (1.15mm,5.09mm) -- (1.24mm,5.18mm) -- (1.34mm,5.27mm) -- (1.45mm,5.38mm) -- (1.57mm,5.49mm) -- (1.69mm,5.62mm) -- (1.82mm,5.75mm) -- (1.96mm,5.91mm) -- (2.12mm,6.09mm) -- (2.29mm,6.29mm) -- (2.49mm,6.53mm) -- (2.71mm,6.81mm) -- (2.98mm,7.16mm) -- (3.30mm,7.62mm) -- (3.75mm,8.30mm) -- (4.87mm,10.35mm) -- cycle;
\draw[tpc-state,line width=1.0pt] (4.87mm,10.35mm) -- (5.89mm,13.05mm) -- (6.21mm,14.34mm) -- (6.40mm,15.40mm) -- (6.50mm,16.32mm) -- (6.53mm,17.14mm) -- (6.52mm,17.89mm) -- (6.47mm,18.59mm) -- (6.38mm,19.22mm) -- (6.27mm,19.81mm) -- (6.12mm,20.36mm) -- (5.95mm,20.86mm) -- (5.76mm,21.32mm) -- (5.54mm,21.74mm) -- (5.31mm,22.12mm) -- (5.06mm,22.47mm) -- (4.79mm,22.77mm) -- (4.51mm,23.04mm) -- (4.21mm,23.28mm) -- (3.90mm,23.48mm) -- (3.58mm,23.64mm) -- (3.25mm,23.77mm) -- (2.91mm,23.87mm) -- (2.57mm,23.93mm) -- (2.21mm,23.96mm) -- (1.85mm,23.95mm) -- (1.49mm,23.91mm) -- (1.12mm,23.84mm) -- (0.75mm,23.73mm) -- (0.37mm,23.60mm) -- (0.00mm,23.43mm) -- (-0.37mm,23.23mm) -- (-0.75mm,23.00mm) -- (-1.12mm,22.75mm) -- (-1.49mm,22.46mm) -- (-1.85mm,22.14mm) -- (-2.21mm,21.80mm) -- (-2.57mm,21.42mm) -- (-2.91mm,21.02mm) -- (-3.25mm,20.60mm) -- (-3.58mm,20.14mm) -- (-3.90mm,19.67mm) -- (-4.21mm,19.17mm) -- (-4.51mm,18.64mm) -- (-4.79mm,18.10mm) -- (-5.06mm,17.53mm) -- (-5.31mm,16.94mm) -- (-5.54mm,16.33mm) -- (-5.76mm,15.70mm) -- (-5.95mm,15.05mm) -- (-6.12mm,14.38mm) -- (-6.27mm,13.69mm) -- (-6.38mm,12.99mm) -- (-6.47mm,12.26mm) -- (-6.52mm,11.52mm) -- (-6.53mm,10.76mm) -- (-6.50mm,9.97mm) -- (-6.40mm,9.15mm) -- (-6.21mm,8.27mm) -- (-5.89mm,7.29mm) -- (-4.87mm,5.60mm) -- (-4.87mm,5.60mm) -- (-3.75mm,4.64mm) -- (-3.30mm,4.39mm) -- (-2.98mm,4.25mm) -- (-2.71mm,4.16mm) -- (-2.49mm,4.10mm) -- (-2.29mm,4.05mm) -- (-2.12mm,4.02mm) -- (-1.96mm,3.99mm) -- (-1.82mm,3.98mm) -- (-1.69mm,3.97mm) -- (-1.57mm,3.96mm) -- (-1.45mm,3.96mm) -- (-1.34mm,3.96mm) -- (-1.24mm,3.97mm) -- (-1.15mm,3.97mm) -- (-1.05mm,3.98mm) -- (-0.96mm,3.99mm) -- (-0.88mm,4.01mm) -- (-0.80mm,4.02mm) -- (-0.72mm,4.04mm) -- (-0.64mm,4.06mm) -- (-0.56mm,4.08mm) -- (-0.49mm,4.10mm) -- (-0.42mm,4.12mm) -- (-0.35mm,4.15mm) -- (-0.28mm,4.17mm) -- (-0.21mm,4.20mm) -- (-0.14mm,4.23mm) -- (-0.07mm,4.26mm) -- (0.00mm,4.29mm) -- (0.07mm,4.33mm) -- (0.14mm,4.36mm) -- (0.21mm,4.40mm) -- (0.28mm,4.44mm) -- (0.35mm,4.48mm) -- (0.42mm,4.53mm) -- (0.49mm,4.58mm) -- (0.56mm,4.63mm) -- (0.64mm,4.68mm) -- (0.72mm,4.74mm) -- (0.80mm,4.80mm) -- (0.88mm,4.87mm) -- (0.96mm,4.94mm) -- (1.05mm,5.01mm) -- (1.15mm,5.09mm) -- (1.24mm,5.18mm) -- (1.34mm,5.27mm) -- (1.45mm,5.38mm) -- (1.57mm,5.49mm) -- (1.69mm,5.62mm) -- (1.82mm,5.75mm) -- (1.96mm,5.91mm) -- (2.12mm,6.09mm) -- (2.29mm,6.29mm) -- (2.49mm,6.53mm) -- (2.71mm,6.81mm) -- (2.98mm,7.16mm) -- (3.30mm,7.62mm) -- (3.75mm,8.30mm) -- (4.87mm,10.35mm) -- cycle;
\fill[tpc-ink] (0.00mm,0.00mm) circle (0.6pt);
\draw[tpc-ink,line width=0.3pt] (6.53mm,17.14mm) -- (15.53mm,23.14mm) node[anchor=west,text=tpc-state,inner sep=1pt] {$\lvert\psi_{210}\rvert^2=\tfrac12\max$};
\draw[tpc-ink,line width=0.3pt] (30.08mm,1.81mm) -- (35.08mm,-2.19mm) node[anchor=west,text=tpc-ink,inner sep=1pt] {Knotenebene $z=0$};
\draw[tpc-ink,line width=0.5pt,-{Stealth[length=3pt]}] (-37.60mm,-28.20mm) -- (-44.49mm,-31.57mm) node[pos=1.25,inner sep=0pt,font=\footnotesize] {$x$};
\draw[tpc-ink,line width=0.5pt,-{Stealth[length=3pt]}] (-37.60mm,-28.20mm) -- (-27.77mm,-30.56mm) node[pos=1.25,inner sep=0pt,font=\footnotesize] {$y$};
\draw[tpc-ink,line width=0.5pt,-{Stealth[length=3pt]}] (-37.60mm,-28.20mm) -- (-37.60mm,-16.93mm) node[pos=1.25,inner sep=0pt,font=\footnotesize] {$z$};
\end{tikzpicture}
\\[2mm]
\parbox{112mm}{\fontsize{9}{11.5}\selectfont\textbf{Abb. 16.2}\enspace
\emph{Die Fläche halber Maximaldichte umschließt nur \(\SI{18}{\percent}\) der Wahrscheinlichkeit.}
Schnitt \(y=0\) durch \(\lvert\psi_{210}\rvert^2\); Kamera V1 (Azimut \ang{35}, Elevation \ang{20}, orthografisch);
\(1\,a_0=\SI{6}{\milli\metre}\). Die Konturlinie ist aus der Formel berechnet.}
\end{center}%
\live{F-16.2}

\noindent Mit \(V\) gleich dem Inneren dieser Fläche ergibt die Bornsche Regel \(p(V)\approx\num{0.177}\);
mehr als vier Fünftel der Wahrscheinlichkeit liegen in der flachen Dichte außerhalb.

\begin{zoomin}{16.V1}{T-16.2b}
{\scshape Vertiefung 16.V1} (warum \(\chi_V(\hat{\mathbf x})\) ein Effekt ist). Die Multiplikation mit der Indikatorfunktion
\(\chi_V\) ist ein Orthogonalprojektor: \(\chi_V^2=\chi_V=\overline{\chi_V}\). Also gilt \(0\le E_V\le\mathbb{1}\),
und \(E_V+E_{\mathbb{R}^3\setminus V}=\mathbb{1}\) ist ein Messgerät mit zwei Ausgängen.
Die Spur reduziert sich für \(\rho=\lvert\psi\rangle\langle\psi\rvert\) auf
\(\langle\psi\vert E_V\psi\rangle=\int_V\lvert\psi\rvert^2\,\dd^3x\). \emph{Zurück:} \(p(V)\) ist das Integral der Dichte über \(V\).
\end{zoomin}
\end{document}
